
\section{理解包含再循环的二维通道流中的传热}
Understanding heat transfer in 2D channel flows including recirculation

\subsection{涡与涡运动的基本概念}
涡流：流体中\textbf{局部}具有\textbf{旋转性质的运动区域}，通常表现为较强涡量集中，流线可能形成局部绕流旋转结构

涡量：速度的旋度定义为流场的涡量，涡量为零称流动无旋，反之有旋。

涡线、涡面、涡管：涡线是处处与涡量方向一致的曲线，多个相邻涡线围成的管状结构称为涡管，垂直于涡线并包含多个涡线的光滑曲面称为涡面。涡量线不穿过涡管的侧面。

涡通量：涡量穿透某曲面的程度，即涡量通过该曲面的总量

涡管强度：某涡管截面的瞬时涡通量

速度环量：流场中速度沿某封闭曲线的线积分，一般取逆时针方向为正。如果某个区域中涡量为零，那么绕其边界的速度环量也为零 → 无旋流；
如果某区域有非零涡通量，则沿其边界的速度环量就不为零 → 有旋流。

涡管强度守恒定理：同一涡管截面上，涡通量相同，涡量场是一个无源场。截面越小涡量越大，流体旋转角速度越大。涡管不能在流体之中产生或终止，只能在流体中形成环形涡环、或始于终于边界面、或延伸至无穷远。
\subsection{文献叙述}
研究对象：包含回流区的二维湍流流动。

SC变换存在未知参数，这些参数实际代表了速度场的几何特征。这些特征将流场描述为channel flow与recirculation wake的组合。

回流区经常出现在通道内障碍物（如挡板、baffle）后方，是工程中常见的非线性区域

SC 映射将上半平面映射到多边形区域，常用于描述复杂边界形状。
本文使用 SC 映射模拟带回流区的流场，将流场抽象为由“主流通道 + 回流尾迹”组成的复几何结构。

SC 映射中的参数是未知的，但它们有明确的几何意义：描述流场中通道宽度、回流区长度、尾迹形状等。
这些参数体现了速度场的几何特征，是建立简化模型（lumped model）的基础。

不止是模拟数据，还尝试将 SC 模型应用到真实测量的速度场上进行参数反推（calibration）

\subsubsection{回流产生的原因}
无旋流的速度场不满足无滑移边界条件。壁面附近的边界层中，速度梯度大但是几何厚度小，黏性效应非常重要。在某些情况下，边界层会从形成它的壁面上分离。在分离点，微小的涡被输运到主流中，如果雷诺数足够高，会形成一个薄层，称为涡层(vortex sheet)。如果涡层底下的区域是有界的，通过涡层传递的湍流动量就会形成回流。回流涡旋是小涡通过涡量汇聚与再组织，在有约束的空间中演化出的大涡结构。

（关于涡层：涡层与切向速度间断面等价，实际上是有很强剪切速度的薄层，显然此层内点点有旋。涡层可以视为无限涡丝组成，涡丝外流体无旋，因此涡层两侧的流体也是无旋的。

无穷长涡丝感生的是平面运动，此时涡丝在平面上表现为一个点涡，平面上点涡感生的速度场表示为$v=\pm\frac{\Gamma
}{2\pi\rho}\theta_{0}$）


\subsection{定义一些重要概念}
滞止涡 stagnation vortex

回流涡 recirculation vortex

主流通道 main flow

涡层 vortex sheet
\begin{itemize}
    \item 附着涡层 attachment vortex sheet（边界层）
    \item 自由涡层 free vortex sheet
    \item 尾涡层 weak vortex 
\end{itemize}

普通分离（逆压梯度引起，也叫挤压分离）、溢流分离（物体后缘）

流动分离：若涡量通过对流离开物体表面，并且随后形成涡片和涡，则存在分离。涡层和涡同时承载运动学动涡量和静涡量。物体表面的边界层方程在各自的分离位置上都是无效的。

有限长度物体上总会发生分离

分离过程通常涉及两个边界层

\subsubsection{分离流和涡流的流动物理和数学模型}
流动物理和数学模型


\begin{enumerate}
    \item 真实流动：现实中感知到的流动，涉及边界层、涡层、涡流，通过扩散和对流的方式输运涡量【做实验】
    \item 模拟黏性流：无黏流与边界层的区别，边界（剪切）层的运动学动涡含量和静涡含量与无黏外流的大小有关【将势流理论的奇点与边界层、剪切层和涡的特性联系起来】
    \item 势流（laplace方程）：不可压无旋无黏流，无全局力作用在物体上（达朗贝尔悖论）【无库塔条件】
    \item 模拟势流：不可压无旋无黏流，以及以涡奇点形式产生的涡层和涡
    \item 可压缩势流*：可压无旋无黏流，以及以涡奇点形式产生的涡层和涡
    \item Euler流：无黏流，涡量仅在弯曲激波中产生，对流涡量输运
    \item 模拟Euler流：无黏流，以及视为涡奇点的涡层和涡
    \item 离散模拟Euler流：无黏流，以及因数值方法的固有黏性而导致的扩散输运引起的涡层和有限厚度的涡
    \item NS流：大多数完全模拟的层流，原则上也包括转捩流（DNS法），涡量的产生和输运，固有的数值黏性可能在所有的类型的扩散输运中引起误差
    \item RANS、URANS流：雷诺平均NS湍流，还包括非定常雷诺平均NS流，我俩的产生和输运，固有的数值黏性
    \item 尺度分解和混合流：大涡模拟等模拟湍流，涡量的产生和输运，与RANS/URANS耦合，固有的数值黏性
\end{enumerate}

通常假设流动有序（只存在基本特征），次要特征（二次三次分离现象与不稳定性）也隐含包括在内。

单元问题（独立流动问题、单物理特性）-构建（基准问题）-子系统（一体化）-系统

\subsubsection{边界层特性}
边界层的烈度、强度
\begin{itemize}
    \item 烈度：边界层的局部涡含量$|\Omega|$。指分离时边界层的局部涡量
    \item 强度：边界层承载的平均表面切向动量$\rho u^{2}$。边界层的强度与局部涡含量无关。
\end{itemize}
二维流动中，主流方向的边界层通常比次边界层具有更大的强度，次边界层在分离区内，通常具有低涡含量和低强度。该边界层在某种意义上是一个诱导边界层，由主边界层的分离引起。

边界层克服逆压梯度的能力越强，其所承载的表面切向动量就越大。经典的实例是给定逆压梯度下，层流边界层比湍流边界层更早分离，这是因为湍流边界层表面切向动量通量更大。

表面切向动量通量的衡量指标是动量厚度，或者说是动量损失厚度。切向边界层速度剖面越饱满，动量损失越小。当然，湍流边界层切向速度剖面比层流边界层更饱满。

壁面相容性条件：切向速度的二阶导数（壁面切向应力的法向导数） 

二维涡方程
\begin{align}
    u\frac{\partial \omega}{\partial x}+v\frac{\partial \omega}{\partial y}&=\frac{1}{Re}\left(\frac{\partial^{2} \omega}{\partial x^{2}}+\frac{\partial^{2} \omega}{\partial y^{2}}\right)\\
    \omega = \frac{\partial v}{\partial x}-\frac{\partial u}{\partial y}&=-\frac{\partial^{2} \psi}{\partial x^{2}}-\frac{\partial^{2} \psi}{\partial y^{2}}\\
    Re&=\frac{\Omega L^{2}}{\nu}
\end{align}
对称monopoles，非线性项消失

rankine涡模型
\begin{align}
    \psi = \psi_{r}\left(1-\frac{r^{2}}{R^{2}}\right)\\
    \psi = \psi_{r}\left(1-\frac{(x-x_{0})^{2}+(y-y_{0})^{2}}{R^{2}}\right)
\end{align}
内核区 ($r<R$)：固体旋转流，$\omega \neq 0$